\chapter{Field expansion for s-dependent magnetic field}
\label{chapter:splineboris-field-expansion}

The \texttt{SplineBoris} element reconstructs a three-dimensional static
magnetic field from the field components and their transverse derivatives on
the reference axis. This chapter describes the expansion used for that
reconstruction. The derivation follows the method presented in
Ref.~\cite{VanDerSchueren2024}, specialized to the straight reference frame
supported by \texttt{SplineBoris}.

\section{Scalar-potential representation}

Inside a current-free aperture, a static magnetic field satisfies
\[
\nabla\times\bm B=0,
\qquad
\nabla\cdot\bm B=0.
\]
Locally, the first equation allows the field to be written as the gradient of
a scalar potential. We use the sign convention employed by the symbolic field
generator,
\begin{equation}
\bm B=\nabla\Phi.
\label{field-expansion:B-from-phi}
\end{equation}
The opposite convention, \(\bm B=-\nabla\Phi\), is equivalent after changing
the sign of \(\Phi\). The divergence equation then gives Laplace's equation,
\begin{equation}
\nabla^2\Phi=0.
\label{field-expansion:laplace}
\end{equation}

The local coordinates \((x,y,s)\) form a straight Cartesian frame, with \(s\)
along the reference axis. Expand the potential in powers of \(y\),
\begin{equation}
\Phi(x,y,s)=\sum_{n=0}^{\infty}\phi_n(x,s)\frac{y^n}{n!}.
\label{field-expansion:phi-series}
\end{equation}
Substitution into Eq.~\eqref{field-expansion:laplace} gives
\begin{equation}
\phi_{n+2}(x,s)
=-\left(\partial_x^2+\partial_s^2\right)\phi_n(x,s).
\label{field-expansion:recurrence}
\end{equation}
Therefore only \(\phi_0\) and \(\phi_1\) are independent; every higher
coefficient follows from the recurrence. For example,
\begin{align}
\Phi={}&\phi_0+y\phi_1
-\frac{y^2}{2!}\left(\partial_x^2+\partial_s^2\right)\phi_0
-\frac{y^3}{3!}\left(\partial_x^2+\partial_s^2\right)\phi_1
+\cdots .
\label{field-expansion:first-terms}
\end{align}
This form also makes clear that the even powers of \(y\) are determined by
\(\phi_0\), while the odd powers are determined by \(\phi_1\).

\section{Independent longitudinal functions}

The two independent functions are expanded in \(x\) as
\begin{align}
\phi_0(x,s)
&=c_0(s)+\sum_{m=1}^{\infty}a_m(s)\frac{x^m}{m!},
\label{field-expansion:phi0}\\
\phi_1(x,s)
&=\sum_{m=1}^{\infty}b_m(s)\frac{x^{m-1}}{(m-1)!},
\label{field-expansion:phi1}
\end{align}
where \(c_0'(s)=b_s(s)\). Applying
Eq.~\eqref{field-expansion:B-from-phi} on the plane \(y=0\) gives
\begin{align}
B_x(x,0,s)
&=\sum_{m=1}^{\infty}a_m(s)\frac{x^{m-1}}{(m-1)!},\\
B_y(x,0,s)
&=\sum_{m=1}^{\infty}b_m(s)\frac{x^{m-1}}{(m-1)!},\\
B_s(0,0,s)&=b_s(s).
\end{align}
Consequently, the longitudinal input functions have a direct interpretation:
\begin{align}
a_m(s)
&=\left.\frac{\partial^{m-1}B_x}{\partial x^{m-1}}
\right|_{x=y=0},
\label{field-expansion:am}\\
b_m(s)
&=\left.\frac{\partial^{m-1}B_y}{\partial x^{m-1}}
\right|_{x=y=0},
\label{field-expansion:bm}\\
b_s(s)&=B_s(0,0,s).
\label{field-expansion:bs}
\end{align}
Thus, knowledge of \(B_s\) on axis and of the successive transverse
derivatives of \(B_x\) and \(B_y\) on axis is sufficient to reconstruct all
three field components around the reference axis. Since the reconstructed
field is obtained from a harmonic scalar potential, it satisfies both static
Maxwell equations used above.
